% ---
% Reemplaza líneas 1191-1229 de Docs/Technical_Report.tex
% Decisiones: N/A (sin cambio)
% ---

\subsection{Circuitos cuánticos variacionales y 
            \textit{quantum embeddings}}

% -----------------------------------------------------------
\subsubsection{Circuitos cuánticos parametrizados (PQC)}
% -----------------------------------------------------------

Un \textit{circuito cuántico parametrizado} (PQC, por sus
siglas en inglés) es un circuito cuántico cuyas compuertas
dependen de un conjunto de parámetros reales
$\boldsymbol{\theta} \in \mathbb{R}^m$ que pueden ajustarse
mediante optimización clásica \cite{cerezo2021variational}.
Un \textit{circuito cuántico variacional} (VQC) es un tipo
específico de PQC que combina un bloque de codificación de
datos $U(\mathbf{x})$ con un bloque variacional entrenable
$W(\boldsymbol{\theta})$, produciendo el estado cuántico:

\begin{equation}
  |\psi(\mathbf{x}, \boldsymbol{\theta})\rangle =
  W(\boldsymbol{\theta})\,U(\mathbf{x})\,
  |0\rangle^{\otimes n}
  \label{eq:vqc_state}
\end{equation}

\noindent donde $n$ es el número de qubits,
$U(\mathbf{x})$ es el \textit{feature map} cuántico que
codifica el vector de características clásico
$\mathbf{x} \in \mathbb{R}^k$ en el estado cuántico inicial,
y $W(\boldsymbol{\theta})$ es el \textit{ansatz} variacional
que introduce correlaciones entre qubits mediante compuertas
entrenables \cite{Schuld2019QuantumML}.

La arquitectura VQC es la más adecuada para el régimen NISQ
actual porque mantiene profundidades de circuito moderadas
y es diferenciable de extremo a extremo mediante el método
de \textit{parameter-shift}, descrito en la
sección~\ref{sec:parameter_shift}.
